\documentclass[10pt,a4paper]{amsart}
\usepackage[margin=19mm]{geometry}
\usepackage{amsmath,amssymb,mathtools}
\usepackage[scr=rsfs,cal=euler]{mathalfa}
\usepackage{xfrac}
\usepackage[unicode,hidelinks,bookmarks=false]{hyperref}
\newcommand{\ie}{i.e.\ }
\newcommand{\cf}{cf.\ }
\newcommand{\resp}{respectively\ }
\newcommand{\Subsection}{\subsection}
\providecommand{\nobreakdash}{\nobreak\hbox{-}}
\setlength{\emergencystretch}{2em}
\multlinegap0pt
\usepackage{mathrsfs}
\usepackage{tikz-cd}
\usepackage{bm}
\usepackage{bbm}
\usepackage{ragged2e}
\usepackage{enumitem}
\setcounter{MaxMatrixCols}{20}
\newtheorem{theorem}{Theorem}
   \newtheorem{corollary}[theorem]{Corollary}
   \newtheorem{lemma}{Lemma}
   \newtheorem*{lemma*}{Lemma} 
   \newtheorem{proposition}[theorem]{Proposition}
   \newtheorem{assumptions}[theorem]{Assumptions}
   \newtheorem{definition}[theorem]{Definition}
   \newtheorem{example}[theorem]{Example}
  \newtheorem{remark}[theorem]{Remark}
\newcommand*{\sbar}{\skew{3}{\bar}}
\newcommand*{\sdot}{\skew{3}{\dot}}
\newcommand{\gv}[1]{\ensuremath{\mbox{\boldmath$ #1 $}}} % for vectors of Greek letters (ancienne définition)
\DeclareMathOperator{\tr}{tr}
\DeclareMathOperator{\rank}{rank}
\DeclareMathOperator{\diag}{diag}
\DeclareMathOperator{\Real}{Re}
\DeclareMathOperator{\Imag}{Im}
\DeclareMathOperator{\Arg}{Arg}
\DeclareMathOperator{\id}{id}
\DeclareMathOperator{\pr}{pr}
\DeclareMathOperator{\SO}{SO}
\DeclareMathOperator{\SL}{SL}
\DeclareMathOperator{\PSL}{PSL}
\DeclareMathOperator{\GL}{GL}
\DeclareMathOperator{\PGL}{PGL}
\DeclareMathOperator{\MCG}{MCG}
\DeclareMathOperator{\SU}{SU}
\DeclareMathOperator{\PSU}{PSU}
\DeclareMathOperator{\Aut}{Aut}
\DeclareMathOperator{\cotan}{cotan}
\DeclareMathOperator{\sech}{sech}
\DeclareMathOperator{\csch}{csch}
\DeclareMathOperator{\arcsec}{arcsec}
\DeclareMathOperator{\arccotan}{arccotan}
\DeclareMathOperator{\arccosec}{arccosec}
\DeclareMathOperator{\arccosh}{arccosh}
\DeclareMathOperator{\arcsinh}{arcsinh}
\DeclareMathOperator{\arctanh}{arctanh}
\DeclareMathOperator{\arcsech}{arcsech}
\DeclareMathOperator{\arccosech}{arccosech}
\DeclareMathOperator{\arccotanh}{arccotanh} 
\newcommand{\dif}{\mathop{}\!\mathrm{d}}
\newcommand{\Dif}{\mathop{}\!\mathrm{D}}
\newcommand\pd[3][]{\frac{\partial\spx{#1}#2}{\partial#3\spx{#1}}}
\newcommand\tpd[3][]{\tfrac{\partial\spx{#1}#2}{\partial#3\spx{#1}}}
\newcommand\dpd[3][]{\dfrac{\partial\spx{#1}#2}{\partial#3\spx{#1}}}
\newcommand{\md}[6]{\frac{\partial\spx{#2}#1}{\partial#3\spx{#4}\partial#5\spx{#6}}}
\newcommand{\tmd}[6]{\tfrac{\partial\spx{#2}#1}{\partial#3\spx{#4}\partial#5\spx{#6}}}
\newcommand{\dmd}[6]{\dfrac{\partial\spx{#2}#1}{\partial#3\spx{#4}\partial#5\spx{#6}}}
\newcommand{\od}[3][]{\frac{\dif\spx{#1}#2}{\dif#3\spx{#1}}}
\newcommand{\tod}[3][]{\tfrac{\dif\spx{#1}#2}{\dif#3\spx{#1}}}
\newcommand{\dod}[3][]{\dfrac{\dif\spx{#1}#2}{\dif#3\spx{#1}}}
\newcommand{\del}[2][-1]{\genericdel(){#1}{#2}}
\newcommand{\set}[2][-1]{\genericdel\{\}{#1}{#2}}
\newcommand{\sbr}[2][-1]{\genericdel[]{#1}{#2}}
\newcommand{\intoc}[2][-1]{\genericdel(]{#1}{#2}}
\newcommand{\intco}[2][-1]{\genericdel[){#1}{#2}}
\newcommand{\eval}[2][-1]{\genericdel.|{#1}{#2}}
\newcommand{\envert}[2][-1]{\genericdel||{#1}{#2}}
\newcommand{\sVert}[1][0]{%
  \ifcase#1\relax
  \rvert\or\bigr|\or\Bigr|\or\biggr|\or\Biggr
  \fi
}
\newcommand{\enVert}[2][-1]{\genericdel\|\|{#1}{#2}}
\newcommand{\fullfunction}[5]{%
  \begin{array}{@{}r@{}r@{}c@{}l@{}}
  #1 \colon & #2 & {}\longrightarrow{} & #3 \\
            & #4 & {}\longmapsto{}     & #5
  \end{array}
}
\begin{document}
\title[Operator-theoretic approach to the partial integration of ra]{Operator-theoretic approach to the partial integration of randomly coupled phase oscillators}
\author{Benjamin Claveau Département de physique, de génie physique et d'optique, Université Laval, Québec (Qc), Canada; Vincent Thibeault, Benjamin Claveau, Antoine Allard,; Patrick Desrosiers}
\date{}
\maketitle
\noindent\emph{Source and licence.} Operator-theoretic approach to the partial integration of randomly coupled phase oscillators. Version arXiv:2607.02617v1 (CC BY 4.0). This material is distributed under the Creative Commons Attribution 4.0 International licence, http://creativecommons.org/licenses/by/4.0/, as recorded in the arXiv OAI licence metadata for this exact version and shown on the abstract page https://arxiv.org/abs/2607.02617v1. Full rights notice: licenses/arxiv-2607.02617-CC-BY-4.0.txt.\par\smallskip
\noindent\textit{Editorial scope.} Counted unit: the original \texttt{lemma} environment \texttt{lem:correspondance\_nonautonome} at source lines 774--789 together with its complete proof at lines 790--797; the delivered document keeps source lines 757--803 so that every referenced label and every citation resolves inside it. Native editable source is the paper's own concatenated TeX; the AMS wrapper is editorial formatting only. No publisher class, upstream script or unrelated artwork is redistributed.
\par\smallskip\noindent\textit{Reference note.} This article ships no compiled bibliography (biblatex); the entries delivered here are composed from the author's own BibTeX (.bib) fields, with no field invented and no entry dropped.
\section{From the Kuramoto model to Riccati equations}
\label{subsec:descriptions}

In Ref.~\cite{Thibeault2026}, we have shown the following lemma.
\begin{lemma}\label{lem:correspondance_complexe_A}
    With $z_j(t) = e^{i\theta_j(t)}$, the initial value problem of the Kuramoto model is equivalent to 
\begin{align}
    \dot{z}_j(t) &= \sum_{k=1}^NA_{jk}z_k(t) - \left(\sum_{k=1}^N\bar{A}_{jk}\bar{z}_k(t) \right)z_j(t)^2 \label{eq:zkurA}\\
    z_j(0) &= e^{i\vartheta_j}  \in \mathbb{T}\label{eq:z0kurA}\,,
\end{align}
where $A$ is a complex matrix of interactions satisfying
\begin{align}
    A = \frac{1}{2}\left(W \circ e^{-i\alpha} + i\diag(\bm{\omega})\right)\,,
\end{align}
where $e^{-i\alpha} = (e^{-i\alpha_{jk}})_{j,k}$, $\bm{\omega} = (\omega_1,...,\omega_N)$, $\circ$ is the element-wise product and $\diag(W) = \diag(\alpha) = \bm{0}$. There exists a constant $a > 0$ such that the problem~(\ref{eq:zkurA}-\ref{eq:z0kurA}) possesses a unique solution $z_1(t),..., z_N(t)$ on $t\in[-a, a]$.
\end{lemma}
To perform the partial integration of the Kuramoto model, it is useful to establish a correspondence with one solution of the model (thus, for a given initial condition) and one solution of a set of Riccati equations representing a non-autonomous dynamical system.

\begin{lemma}\label{lem:correspondance_nonautonome}
    Let the initial value problem of Riccati be
    \begin{align}
        \dot{u}_j(t) &= q_j(t) - \overline{q_j(t)}u_j(t)^2\,,\qquad j\in\{1,...,N\}\label{eq:riccati}\\
        u_j(0) &= d_j \in \mathbb{T} \label{eq:uCI}
    \end{align} % $q_{-1}(t) = \overline{q_1(t)}$
    such that $u_j(t)\in\mathbb{T}$ for all $t \in [-a, a]$ with $a > 0$ and the arbitrary complex function $q_j$ defines the problem. Moreover, consider some solution of the Kuramoto model [Lemma~\ref{lem:correspondance_complexe_A}] $z_1(t), ..., z_N(t)$ for $t \in [-a, a]$ with $z_1(0) = \exp(i\vartheta_1), ..., z_N(0) = \exp(i\vartheta_N)$ and a complex matrix of interactions $A$. If the Riccati problem in Eqs.~(\ref{eq:riccati}-\ref{eq:uCI}) satisfies $d_j = \exp(i\vartheta_j)$ and is defined with
    \begin{align}\label{eq:pj_to_qj}
        q_j(t) = \textstyle{\sum_{k=1}^N} A_{jk}z_k(t)\,, \qquad j\in\{1,...,N\}\,,%\qquad\text{and}\qquad q_{-1}(t) = \overline{q_1(t)}\,.
    \end{align}
    then its unique solution coincides with the one for the Kuramoto model, i.e.,
    \begin{align*}
        u_j(t) = z_j(t)\,,\qquad j\in\{1,...,N\}\,,
    \end{align*}
    on $t \in [-a, a]$. 
\end{lemma}

\begin{proof}
    Substituting $d_j$, $q_j$ and $\bar{q}_j$ in Eqs.~(\ref{eq:riccati}-\ref{eq:uCI}) gives
    \begin{align}
        \dot{u}_j(t) &= \sum_{k=1}^N A_{jk}z_k(t)  - \sum_{k=1}^N \bar{A}_{jk}\bar{z}_k(t)u_j(t)^2\,, \qquad j\in\{1,...,N\}\label{eq:u1ric}\\
        u_j(0) &= \exp(i\vartheta_j) = z_j(0)\label{eq:u2ric}\,.
    \end{align} 
    Clearly, $u_j(t) = z_j(t)$ for all $j$ and $t$, because $(u_j(t))_j$ is a solution of Eqs.~(\ref{eq:u1ric}-\ref{eq:u2ric}), which are equivalent to Eqs.~(\ref{eq:zkurA}-\ref{eq:z0kurA}) [Lemma~\ref{lem:correspondance_complexe_A}]. Since $f_j(x_j, t) = q_j(t) - \overline{q_j(t)}x_j^2$ for each $j\in\{1, ..., N\}$ is continuous in $t$ and continuously differentiable in $x_j$ for all $x_j\in\mathbb{T}$, then the fundamental existence-uniqueness theorem for non-autonomous systems~\cite[p.77]{Perko2001} guarantees the uniqueness of the solution (recall from Lemma~\ref{lem:correspondance_complexe_A} that $z_1(t)$,...,$z_N(t)$ is also unique).
\end{proof}

\begin{remark}
    One must notice that if either the condition $d_j = \exp(i\vartheta_j)$ or $q_j(t) = \textstyle{\sum_{k=1}^N} A_{jk}z_k(t)$ is not satisfied, then no correspondence exists between the Riccati problem and the Kuramoto model. In perhaps more intuitive terms, different trajectories of the Kuramoto dynamics are related to different Riccati equations (i.e., $q_1(t)$,...,$q_N(t)$ are different for each Kuramoto solution) under the condition that the initial conditions coincide.
\end{remark}
% \begin{lemma}
%     If $d_j \neq \exp{i\vartheta_j}$ for all $j$ while the Riccati problem is still defined with Eq.~\eqref{eq:pj_to_qj}, the unique solution $u_j(t)$ will 
% \end{lemma}

\begin{thebibliography}{99}
\bibitem[]{Perko2001} Perko, L.. Differential Equations and Dynamical Systems. (2001)
\bibitem[]{Thibeault2026} Thibeault, V.; Claveau, B.; Allard, A.; Desrosiers, P.. Kuramoto meets Koopman: Constants of motion, symmetries, and network motifs. arXiv:2504.06248 (2026)
\end{thebibliography}
\end{document}
